\section{Method}
\label{sec:method}
We segment the region of a generated image that belongs to a class word $w$ (\eg ``sheep''). All attention is recorded while the image is generated; nothing is trained.

\paragraph{Attention in the denoiser.}
At denoising call $t$, an attention module $m$ with $n_m{=}s_m^2$ image positions forms per-head queries $Q$ and keys $K$. Cross-attention uses text keys from an attribution prompt that contains $w$~\cite{marcosmanchon2024ovam}, and self-attention uses image keys:
\begin{align}
a_{t,m}&=\tfrac1{8}\textstyle\sum_{h}\operatorname{softmax}\!\big(Q_hK_h^\top/\sqrt d\big)_{[:,\,w]}\in\mathbb R^{n_m},\label{eq:attention}\\
A_{t,h}^{s}&=\operatorname{softmax}\!\big(Q_hK_h^\top/\sqrt d\big)\in\mathbb R^{s^2\times s^2}.\nonumber
\end{align}
Here $[:,w]$ averages the columns of all tokenizer pieces of $w$ (softmax over the 77 padded text positions, conditional branch). $A^{s}_{t,h}$ is row-stochastic, so $A^s_{t,h}x$ replaces each value of a map $x$ by an affinity-weighted average. We write $\bar A^s_t=\tfrac18\sum_h A^s_{t,h}$ for the head mean, $U$ for bilinear $2\times$ upsampling and $D$ for area downsampling.

\subsection{Semantic seed}
\label{sec:method_seed}
CA modules exist at 8, 16, 32 and 64 pixels. We use the eleven modules at 8, 16 and 64 pixels, $\mathcal M$ (Sec.~\ref{sec:seed}), and all 31 calls, and weight them per call:
\begin{align}
c&=\operatorname{norm}\Big(\tfrac1{31}\textstyle\sum_{t=0}^{30}\sum_{m\in\mathcal M}w_{t,m}\,U_{64}(a_{t,m})\Big),\nonumber\\
w_{t,m}&=\big[\langle \tilde S_{t,m},P_t\rangle\big]_+\big/\textstyle\sum_{m'\in\mathcal M}\big[\langle \tilde S_{t,m'},P_t\rangle\big]_+,
\label{eq:a3}
\end{align}
where $\operatorname{norm}$ is min--max normalization and $U_{64}$ resizes to $64\times64$. $\tilde S_{t,m}$ is the head-mean SA affinity of the transformer block containing CA module $m$, resized to the reference grid. $P_t=\operatorname{softmax}(FF^\top/\sqrt d)$ is a pseudo-affinity of the decoder's 32-pixel query features $F$ at the same call. These ``A3'' weights, a GoCA-style adaptation~\cite{meng2026making}, favour CA modules whose spatial grouping matches the network's own. We also tested GoCA's transposed orientation; it changes the seed by less than $10^{-3}$ (Pearson $r{=}0.99999$) and the results not at all (supplementary material).

\subsection{Coarse-to-fine propagation with a one-sided cap}
\label{sec:method_c2f}
A single 64-pixel SA module propagates the seed with fine detail but also leaks along fine textures (Sec.~\ref{sec:envelope}). We therefore start at the coarsest SA level that still separates objects, and pass the map up one resolution at a time. For each SA call $t\in\mathcal T=\{10,\dots,20\}$, with $c_{16}=D(c)$,
\begin{equation}
r_t=\bar A^{16}_t c_{16},\qquad q_t=\bar A^{32}_t\,U(r_t),\qquad z_t=\operatorname{cap}(r_t,q_t),
\label{eq:c2f}
\end{equation}
and the final map is $y=\tfrac1{|\mathcal T|}\sum_t \bar A^{64}_t\,U(z_t)$. The cap compares every $2\times2$ group $g$ of 32-pixel children with its 16-pixel parent:
\begin{equation}
\operatorname{cap}(r,q)_k=q_k\cdot\min\!\Big(1,\ \frac{r_{g(k)}}{\tfrac14\sum_{k'\in g(k)}q_{k'}+\epsilon}\Big),
\label{eq:cap}
\end{equation}
with $\epsilon=10^{-12}$. Bilinear upsampling preserves local means, so $4r_g$ is the parent's mass at child resolution. Equation~\eqref{eq:cap} is then the generalized-KL projection of $q$ onto the \emph{upper} constraints $\sum_{k\in g}z_k\le 4r_g$: groups within budget are untouched, and groups over budget are scaled down uniformly. The clipped mass is discarded, not redistributed. Two alternatives frame this choice. Without a cap, the 32-pixel level may place evidence wherever its affinities lead, including into background that the coarse level had excluded. With \emph{exact conservation}, $\sum_{k\in g}z_k=4r_g$, the children must also fill every parent's budget, including the budget of a coarse parent that was wrong. Section~\ref{sec:cap} shows that only the one-sided version helps. We cap the $16\to32$ transition only, and leave $32\to64$ uncapped so that fine SA can refine boundaries.

\subsection{Mask and boundary refinement}
\label{sec:method_crf}
The map is min--max normalized, bilinearly upsampled to image size and thresholded at a global $\tau$ chosen on VOC-Sim training folds ($\tau{=}0.55$). DenseCRF~\cite{krahenbuhl2011crf} then refines the binary mask $b$ by minimizing
\begin{align}
E(x)={}&\textstyle\sum_i\psi_u(x_i\mid b_i)+\sum_{i<j}[x_i\ne x_j]\,k(i,j),\label{eq:crf}\\
k(i,j)={}&w_1\exp\!\Big(-\tfrac{\|p_i-p_j\|^2}{2\theta_\alpha^2}-\tfrac{\|I_i-I_j\|^2}{2\theta_\beta^2}\Big)\nonumber\\
&+w_2\exp\!\Big(-\tfrac{\|p_i-p_j\|^2}{2\theta_\gamma^2}\Big),\nonumber
\end{align}
with positions $p$, colours $I$ and OVAM's fixed parameters ($w_1{=}10$, $\theta_\alpha{=}80$, $\theta_\beta{=}13$, $w_2{=}3$, $\theta_\gamma{=}3$, five mean-field iterations)~\cite{marcosmanchon2024ovam}. We call this step R2. No ground truth enters extraction, propagation or refinement.
